% !TEX program = pdflatex
% The Fock commutator of the CT generator, in the reader's own order of steps
% and Fock convention: [F,A] = [F,T] + h.c., the two blocks, and the collection
% into C^q_p and C^{rs}_{pq} by enlarging the summation range and restricting
% it again with deltas.
\documentclass[11pt]{article}
\usepackage[T1]{fontenc}
\usepackage[utf8]{inputenc}
\usepackage{lmodern}
\usepackage{amsmath,amssymb}
\usepackage[margin=1in]{geometry}
\usepackage{booktabs}
\usepackage[hidelinks]{hyperref}

\newcommand{\E}[2]{\hat E^{#1}_{#2}}
\newcommand{\hc}{\mathrm{h.c.}}
\newcommand{\occ}{i}
\newcommand{\vir}{a}

\title{\bfseries The Fock commutator of the CT generator\\[2pt]
       \large $[\hat F,\hat A]$, its two blocks, and the coefficients $C^q_p$,
       $C^{rs}_{pq}$}
\author{}
\date{}

\begin{document}
\maketitle
\vspace{-2.5em}

\noindent
This is the derivation as it is written out by hand, with the algebra checked
against explicit matrices. The convention is the one used there and in
\texttt{formula.tex}: the Fock tensor carries its indices in the order
$f^p_q$, the generator carries the \emph{excitation} operator, and
\begin{equation}
  \hat F=f^p_q\,\E{q}{p},
  \qquad
  \hat T_1=t^\vir_\occ\,\E{\vir}{\occ},
  \qquad
  \hat T_2=\tfrac12\,t^{\vir\beta}_{\occ j}\,\E{\vir\beta}{\occ j},
  \label{eq:ops}
\end{equation}
with $\occ,j$ occupied, $\vir,\beta$ virtual and $t^{\vir\beta}_{\occ j}$
antisymmetric in $(\vir,\beta)$ and in $(\occ,j)$.
$\hat A=\hat T-\hat T^\dagger$ is anti-Hermitian and $\hat T=\hat T_1+\hat T_2$.

\medskip
\noindent
\textbf{Index ranges.} Throughout, the \emph{general} indices $p,q,r,s$ run
over the \emph{whole} space, which is the union of the occupied and the virtual
sets,
\begin{equation}
  p,q,r,s\in\{\occ\}\cup\{\vir\},
  \label{eq:ranges}
\end{equation}
so a general index is an occupied one or a virtual one and nothing else. The
$\hat E$ operators in the general basis, $\E{q}{p}$, therefore carry four
families of blocks --- occupied--occupied, occupied--virtual,
virtual--occupied and virtual--virtual --- and a Kronecker delta is what selects
one of them. This is the convention that makes the enlargement of \S4 work: the
general range contains both the occupied and the virtual parts, so a restricted
sum can be extended to it and then cut back with a delta.

\section*{1. The reduction}

\begin{equation}
  [\hat F,\hat A]=[\hat F,\hat T]+[\hat F,\hat T]^\dagger
   =[\hat F,\hat T_1]+[\hat F,\hat T_2]+\hc,
  \label{eq:reduce}
\end{equation}
because $\hat A=\hat T-\hat T^\dagger$ and
$([\hat F,\hat T])^\dagger=-[\hat F,\hat T^\dagger]$. Only the two $\hat T$
blocks have to be computed; the other half of $\hat C=[\hat F,\hat A]$ is their
conjugate.

\section*{2. \texorpdfstring{$[\hat F,\hat T_1]$}{[F,T1]}}

Writing the operators out,
\begin{equation}
  [\hat F,\hat T_1]
   =\bigl[f^p_q\E{q}{p},\,t^\vir_\occ\E{\vir}{\occ}\bigr]
   =f^p_q\,t^\vir_\occ\bigl[\E{q}{p},\E{\vir}{\occ}\bigr],
  \label{eq:t1setup}
\end{equation}
and with the $(1,1)$ rule
\begin{equation}
  \bigl[\E{q}{p},\E{s}{r}\bigr]
   =\delta^s_p\,\E{q}{r}-\delta^q_r\,\E{s}{p},
  \label{eq:c11}
\end{equation}
\begin{equation}
  \bigl[\E{q}{p},\E{\vir}{\occ}\bigr]
   =\delta^\vir_p\,\E{q}{\occ}-\delta^q_\occ\,\E{\vir}{p}.
  \label{eq:c11apply}
\end{equation}
The two deltas contract one index of $f$ in each term --- in the first the
\emph{lower} index of $f$ becomes the virtual one, in the second the
\emph{upper} index becomes the occupied one --- so
\begin{equation}
  \boxed{\;[\hat F,\hat T_1]
   =t^\vir_\occ\bigl(f^\vir_q\,\E{q}{\occ}-f^p_\occ\,\E{\vir}{p}\bigr).\;}
  \label{eq:FT1}
\end{equation}

\noindent\textbf{Canonical Fock.} With $f^p_q=\epsilon_p\delta^p_q$,
\begin{equation}
  [\hat F,\hat T_1]=t^\vir_\occ\,(\epsilon_\vir-\epsilon_\occ)\,\E{\vir}{\occ},
  \label{eq:FT1canon}
\end{equation}
and the condition on the occupied--virtual block,
$f^\occ_\vir+[\hat F,\hat T_1]^\occ_\vir=0$, gives the singles amplitude
\begin{equation}
  t^\vir_\occ=\frac{f^\occ_\vir}{\epsilon_\occ-\epsilon_\vir},
\end{equation}
the amplitude of the handwritten singles note.

\section*{3. \texorpdfstring{$[\hat F,\hat T_2]$}{[F,T2]}}

In the same way,
\begin{equation}
  [\hat F,\hat T_2]
   =\tfrac12\,f^p_q\,t^{\vir\beta}_{\occ j}
     \bigl[\E{q}{p},\E{\vir\beta}{\occ j}\bigr],
\end{equation}
the $\tfrac12$ being the one in $\hat T_2$ of \eqref{eq:ops}.
with the $(1,2)$ rule
\begin{equation}
  \bigl[\E{q}{p},\E{tu}{rs}\bigr]
   =\delta^t_p\,\E{qu}{rs}+\delta^u_p\,\E{qt}{sr}
    -\delta^q_r\,\E{tu}{ps}-\delta^q_s\,\E{ut}{pr},
  \label{eq:c12}
\end{equation}
so
\[
  \bigl[\E{q}{p},\E{\vir\beta}{\occ j}\bigr]
   =\delta^\vir_p\,\E{q\beta}{\occ j}+\delta^\beta_p\,\E{q\vir}{j\occ}
    -\delta^q_\occ\,\E{\vir\beta}{pj}-\delta^q_j\,\E{\beta\vir}{p\occ},
\]
so that
\begin{equation}
  [\hat F,\hat T_2]
   =\tfrac12\,f^p_q\,t^{\vir\beta}_{\occ j}
     \Bigl(\delta^\vir_p\,\E{q\beta}{\occ j}+\delta^\beta_p\,\E{q\vir}{j\occ}
      -\delta^q_\occ\,\E{\vir\beta}{pj}-\delta^q_j\,\E{\beta\vir}{p\occ}\Bigr),
  \label{eq:c12apply}
\end{equation}
where the $\tfrac12$ is the one from $\hat T_2$ and must not be dropped: it is
these four terms that are collected in \S4, so the factor stands in front of
the sum. Each delta contracts one index of $f$ with the operator, which gives
the block in the form
\begin{equation}
  \boxed{\;
  \begin{aligned}
    [\hat F,\hat T_2]=\tfrac12\Bigl(
     &+f^\vir_q\,t^{\vir\beta}_{\occ j}\,\E{q\beta}{\occ j}
      +f^\beta_q\,t^{\vir\beta}_{\occ j}\,\E{q\vir}{j\occ}\\
     &-f^\occ_p\,t^{\vir\beta}_{\occ j}\,\E{\vir\beta}{pj}
      -f^j_p\,t^{\vir\beta}_{\occ j}\,\E{\beta\vir}{p\occ}\Bigr).
  \end{aligned}\;}
  \label{eq:FT2}
\end{equation}
The first and the second term are related by the antisymmetry of $t$
(relabelling $\vir\leftrightarrow\beta$ and $\occ\leftrightarrow j$), and so
are the third and the fourth. This relabelling relates the two \emph{sums}; it
is not available term by term, and term by term is how the block is used. At
fixed values of the summation indices the four terms are different operators
with different index patterns --- $\E{q\beta}{ij}$ and $\E{qa}{ji}$ in the first
pair, $\E{\vir\beta}{pj}$ and $\E{\beta\vir}{pi}$ in the second --- so each of
them restricts the general range to a different block, and each carries its own
coefficient. Nothing is cancelled here: cancelling a pair would merge the two
coefficients into one and lose the index patterns that \S4 has to read
$C^{rs}_{pq}$ off.

\noindent\textbf{Canonical Fock.} With $f^p_q=\epsilon_p\delta^p_q$ every one of
the four terms of \eqref{eq:FT2} collapses onto $\E{\vir\beta}{\occ j}$, and the
sum of the four gives
\begin{equation}
  [\hat F,\hat T_2]
   =\tfrac12\,t^{\vir\beta}_{\occ j}
     \bigl(\epsilon_\vir+\epsilon_\beta-\epsilon_\occ-\epsilon_j\bigr)
     \E{\vir\beta}{\occ j}.
  \label{eq:FT2canon}
\end{equation}
Each of the four contributes one of the four orbital energies, which is what
the uncollected form makes visible and what the $\tfrac12$ in front of the sum
is doing there.

\section*{4. Collecting into \texorpdfstring{$C^q_p$}{Cqp} and
\texorpdfstring{$C^{rs}_{pq}$}{Crs pq}}

The result of \S2 and \S3 is a sum of $\hat E$ operators whose indices are
partly restricted --- $\E{q}{\occ}$ has an occupied lower index,
$\E{\vir}{p}$ a virtual upper one. To read off coefficients in the general
basis
\begin{equation}
  [\hat F,\hat A]=C^q_p\,\E{q}{p}+C^{rs}_{pq}\,\E{rs}{pq},
  \label{eq:Cform}
\end{equation}
each restricted sum is enlarged to the general range \eqref{eq:ranges} and
restricted again with a delta. The enlargement is legitimate precisely because
$p,q$ cover the whole space: extending an occupied sum to all of
$\{\occ\}\cup\{\vir\}$ adds the virtual labels, and the delta removes them
again. For the one-body block,
\begin{align}
  \sum_{\occ,\vir,q}t^\vir_\occ f^\vir_q\,\E{q}{\occ}
   &=\sum_{p,q,\occ,\vir}t^\vir_\occ f^\vir_q\,\delta^p_\occ\,\E{q}{p},\notag\\
  -\sum_{\occ,\vir,p}t^\vir_\occ f^p_\occ\,\E{\vir}{p}
   &=-\sum_{p,q,\occ,\vir}t^\vir_\occ f^p_\occ\,\delta^\vir_q\,\E{q}{p},
  \label{eq:enlarge1}
\end{align}
so the $[\hat F,\hat T]$ half of the one-body coefficient is
\begin{equation}
  C^{[T]\,q}_p=\sum_{\occ,\vir}\Bigl(t^\vir_\occ f^\vir_q\,\delta^p_\occ
   -t^\vir_\occ f^p_\occ\,\delta^\vir_q\Bigr).
  \label{eq:C1}
\end{equation}
The conjugate half contributes the same two structures once more with the
operator's indices exchanged, so the complete one-body coefficient is
\begin{equation}
  \boxed{\;C^q_p=
    \underbrace{f^\vir_q\,t^\vir_p}_{p\ \text{occupied}}
   -\underbrace{f^p_\occ\,t^q_\occ}_{q\ \text{virtual}}
   +\underbrace{f^\vir_p\,t^\vir_q}_{q\ \text{occupied}}
   -\underbrace{f^q_\occ\,t^p_\occ}_{p\ \text{virtual}}\;}
  \label{eq:C1result}
\end{equation}
two terms from $[\hat F,\hat T_1]$ and two from its conjugate.

The two-body coefficient is collected the same way, from the four terms of
$[\hat F,\hat T_2]$ in \eqref{eq:FT2} and the four of its conjugate, each on its
own index pattern: $\E{q\beta}{\occ j}$ and $\E{\occ j}{q\beta}$,
$\E{q\vir}{j\occ}$ and $\E{j\occ}{q\vir}$, $\E{\vir\beta}{pj}$ and
$\E{pj}{\vir\beta}$, $\E{\beta\vir}{p\occ}$ and $\E{p\occ}{\beta\vir}$, carrying
the coefficients $\pm\tfrac12 f t$ of that equation. Nothing is cancelled
between them, for the reason given in \S3: the eight are different operators,
and it is their eight patterns that build up $C^{rs}_{pq}$.

With both coefficients complete, the result of the derivation is
\begin{equation}
  \boxed{\;[\hat F,\hat A]=C^q_p\,\E{q}{p}+C^{rs}_{pq}\,\E{rs}{pq}.\;}
  \label{eq:result}
\end{equation}
The two deltas cut the general range back down: $\delta^p_\occ$ keeps the
\emph{occupied} part of it for the lower index, so the first term contributes
only for occupied $p$; $\delta^\vir_q$ keeps the \emph{virtual} part for the
upper index, so the second contributes only for virtual $q$. Every coefficient
below is therefore written over the whole space, with the index ranges carried
by the deltas rather than by the summation limits --- which is the point of the
rewriting.

The same enlargement applies to the four terms of \eqref{eq:c12apply}, each
operator now carrying four restrictions, one per index, which is why
$C^{rs}_{pq}$ comes out as a sum of delta products rather than single deltas.
Two cautionary points, both of which have already bitten this derivation:

\begin{itemize}
  \item \textbf{Keep each delta on the index it restricts.} The enlargement
        fixes the index of $\hat E$, not the index of $f$: the first term of
        \eqref{eq:C1} is $f^\vir_q\delta^p_\occ$, and writing it as
        $f^\occ_q\delta^p_\vir$ instead is wrong while still agreeing for
        symmetric $f$.
  \item \textbf{Canonicalise the two-body operator before reading the
        coefficient.} Once $p,q$ run over all orbitals the same $\E{rs}{pq}$ is
        reached from more than one term; fixing one index ordering first is
        what prevents $C^{rs}_{pq}$ and $C^{sr}_{qp}$ being counted separately.
\end{itemize}

\noindent
These coefficient tensors are the $\bar h$ and $\bar g$ of the CT Hamiltonian
once the conjugate half of \eqref{eq:reduce} is added.

\section*{5. What is verified}

\path{tests/test_F_A_commutator.py} builds $\hat F,\hat T_1,\hat T_2$ as
explicit spin-free matrices on four orbitals ($\{0,1\}$ occupied, $\{2,3\}$
virtual) with random tensors, and compares the closed forms against the matrix
commutators:
\begin{center}
\footnotesize
\begin{tabular}{@{}lrr@{}}
\toprule
quantity & max$|\Delta|$ & relative \\
\midrule
$[\hat F,\hat T_1]$, \eqref{eq:FT1} & $1.1\times10^{-15}$ & $2.6\times10^{-16}$ \\
$[\hat F,\hat T_2]$, \eqref{eq:FT2}, the four terms uncancelled
  & $3.6\times10^{-15}$ & $1.5\times10^{-16}$ \\
the same block with the pairs folded & $3.6\times10^{-15}$ & $1.5\times10^{-16}$ \\
$C^q_p$ of \eqref{eq:C1result}, one-body part & --- & $2.5\times10^{-16}$ \\
\eqref{eq:result} against $[\hat F,\hat A]$ & $4.4\times10^{-15}$ & $6.5\times10^{-16}$ \\
\bottomrule
\end{tabular}
\end{center}

\noindent
The third and fourth rows are the two listings of the same operator: leaving the
four terms uncancelled and folding the pairs agree to the same precision, which
is why the note can keep them separate without changing anything.

\noindent
\textbf{The result is verified.} \path{tests/test_F_A_result.py} takes the two
blocks in closed form, adds the conjugate half, and collects every term into the
general basis by running each term over its own index ranges --- the
enlarge-and-restrict step done mechanically, so that no delta has to be
transcribed by hand --- then reconstructs $C^q_p\E{q}{p}+C^{rs}_{pq}\E{rs}{pq}$
and compares it with the matrix commutator $[\hat F,\hat A]$:
$4.4\times10^{-15}$ absolute, $6.5\times10^{-16}$ relative. The compact
four-term form \eqref{eq:C1result} of the one-body coefficient is checked
separately, $2.5\times10^{-16}$.

One convention point the check turned up, worth stating because it is easy to
overlook: with $\hat F=f^p_q\E{q}{p}$ the operator is Hermitian only when $f$ is
symmetric (real), which is the physical case. The step
$[\hat F,\hat A]=[\hat F,\hat T]+[\hat F,\hat T]^\dagger$ uses that
Hermiticity, and it fails for a non-symmetric $f$ --- the check therefore
symmetrises $f$ and asserts that $\hat F$ is Hermitian before comparing.

\section*{6. The other Fock convention}

The derivation above uses $\hat F=f^p_q\E{q}{p}$. With the aligned order
$\hat F=f^q_p\E{q}{p}=f^q_p\hat a^\dagger_q\hat a_p$, in which
$f^q_p=\langle p|\hat f|q\rangle$, the same derivation gives
\begin{equation}
  [\hat F,\hat T_1]=t^\vir_\occ\bigl(f^p_\vir\E{p}{\occ}-f^\occ_p\E{\vir}{p}\bigr),
  \qquad
  [\hat F,\hat T_2]=t^{\vir\beta}_{\occ j}\bigl(f^p_\vir\E{p\beta}{\occ j}
   -f^\occ_p\E{\vir\beta}{pj}\bigr).
  \label{eq:aligned}
\end{equation}
The two sets are related by $f\to f^{\mathsf T}$, so they agree whenever $f$ is
symmetric --- in particular for the diagonal canonical Fock, which is where a
first check is usually made. That is what makes the distinction easy to lose;
\path{tests/test_F_A_commutator.py} checks both conventions and asserts that
each closed form deviates by a relative amount of order one under the other
one's $\hat F$.

\begin{thebibliography}{9}
\bibitem{Kutzelnigg:1997} W.~Kutzelnigg and D.~Mukherjee,
  \emph{Normal order and extended Wick theorem for a multiconfiguration
  reference wave function}, J.~Chem.~Phys.\ \textbf{107}, 432 (1997).
\bibitem{Yanai:2012cj} T.~Yanai and T.~Shiozaki,
  \emph{Canonical transcorrelated theory with projected Slater-type geminals},
  J.~Chem.~Phys.\ \textbf{136}, 084107 (2012).
\end{thebibliography}

\end{document}
